There is an exact favorable early-kill calculation, but it does not close the feasible-projector theorem.  Conditional on \(A\), put \(X_j=Z_j-1/2\).  For \(A_{ij}=1\), exchangeability of the \(N_i\) neighbor assignments and the already proved identity \(\mathbb E_Z[q_{i,h}V_i\mid A]=1\) imply
\[
 1=\mathbb E_Z[q_{i,h}V_i\mid A]
 =N_i^{-1}\sum_{j:A_{ij}=1}\mathbb E_Z[q_{i,h}X_j\mid A]
 =\mathbb E_Z[q_{i,h}X_j\mid A].
\]
For \(A_{ij}=0\), independence gives \(\mathbb E_Z[q_{i,h}X_j\mid A]=0\).  Since \(X_j\) is a centered two-point variable with variance \(1/4\), the first Hoeffding projection is therefore exactly
\[
 q_{i,h}^{(1)}=4\sum_{j\in\mathcal V_n}A_{ij}X_j,
\]
uniformly in \(h\).  The score is odd under the simultaneous sign reversal of all neighbor assignments.  A conditional expectation depending on two binary coordinates consequently has only odd Walsh monomials; after subtracting its two first-order terms, its second Hoeffding projection is exactly zero.  Hence, for any deterministic projector \(\Pi(A)\),
\[
 \bigl\{(I_n-\Pi(A))q_h\bigr\}^{(1)}
 =4(I_n-\Pi(A))AX,
 \qquad
 \bigl\{(I_n-\Pi(A))q_h\bigr\}^{(2)}=0,
\]
and
\[
 \mathbb E_Z\left[\left\|4(I_n-\Pi(A))AX\right\|_2^2\mid A\right]
 =4\left\|(I_n-\Pi(A))A\right\|_{\mathrm F}^2.
\]
This verifies that the first two projections introduce no bandwidth-dependent term and explains the desired \(d/n\) graph floor after removal of the rank-\(r\) signal.  What remains unproved is a uniform bound for the third-and-higher projections after multiplication by the outcomes, together with the error from replacing the population signal projector by the empirical, leave-one-out projector.  Moreover, \({\rm def:pc-handle}\) presently says only ``use leave-one-edge-out eigenvector linearization'' and does not specify a measurable map or formulas for either certificate.  Consequently there is no defined construction whose bias, variance, running time, or rate can yet be proved.